% Generated draft: requires mathematical and source review.
\documentclass[11pt]{article}
\usepackage[margin=25mm]{geometry}
\usepackage{fontspec}
\setmainfont{Latin Modern Roman}
\usepackage{amsmath,amssymb,mathtools,unicode-math}
\setmathfont{Latin Modern Math}
\usepackage{enumitem,needspace,xurl,hyperref}
\setlength{\parindent}{0pt}
\setlength{\parskip}{5pt}
\setlength{\emergencystretch}{3em}
\newcommand{\N}{\mathbb{N}}
\newcommand{\Z}{\mathbb{Z}}
\newcommand{\Q}{\mathbb{Q}}
\newcommand{\R}{\mathbb{R}}
\newcommand{\C}{\mathbb{C}}
\newcommand{\F}{\mathbb{F}}
\newcommand{\A}{\mathbb{A}}
\newcommand{\PP}{\mathbb{P}}
\newcommand{\E}{\mathbb{E}}
\newcommand{\eps}{\varepsilon}
\newcommand{\dd}{\,\mathrm{d}}
\begin{document}
\section*{UnsolvedMath ID: 30000258}
\textbf{Source problem code:} catalogue entry
\subsection{Definitions and mathematical statement}
The source asks for a forcing theory for stationary curves of a uniformly parabolic PDE. For an evolution equation written abstractly as \(\partial_tu=\mathcal F[u]\), a stationary solution satisfies \(\mathcal F[u]=0\), together with the prescribed boundary or periodicity conditions. In a scalar one-dimensional setting its graph is a curve, but the record does not specify that setting or define stationary curve.

Let \(\Sigma\) denote the given skeleton of stationary curves, in the source's intended sense. The requested task is to determine the number and types of additional stationary curves whose existence follows from the skeleton data and the hypotheses on the PDE. Formally, once an admissible equation class \(\mathcal C_\Sigma\) realizing the same specified skeleton data has been fixed, a forced lower bound would have the form
\[
\forall\mathcal F\in\mathcal C_\Sigma,\qquad
\#\{\text{stationary curves of }\mathcal F\text{ outside }\Sigma\}\geq N(\Sigma),
\]
possibly refined by source-defined topological or geometric classes. This is a template for the requested forcing conclusion, not an assertion of a particular bound or a choice of skeleton equivalence. The actual stationary solution set need not be finite, and the source does not identify curves up to any stated equivalence.

The PDE, domain, boundary conditions, uniform parabolicity assumptions, skeleton definition, preserved data, and meaning of which curves are forced are not reproduced. In particular, geometric realization, ordering, isotopy, or braid-type data must not be chosen silently. The supplied assessment checked 2026-08-21 expressly says that the broad request is not a single fully specified theorem and that no result settling it in the stated generality was located.
\subsection{Short English statement}
Determine what additional stationary curves must exist once a specified skeleton is present in a uniformly parabolic equation. The equation class, skeleton data, and curve equivalence must be defined before numerical or topological forcing conclusions can be stated.
\subsection{Sources}
\begin{itemize}
\item Local UnsolvedMath registry, permanent ID 30000258, catalogue entry, \emph{Stationary-Curve Skeletons in Uniformly Parabolic Equations}.
\item \emph{Topological and Variational Methods for Differential Equations} (2006), Oberwolfach Reports, pp. 1601--1678, DOI \url{https://doi.org/10.4171/owr/2005/29}.
\item Supplied literature triage checked 2026-08-21 cites the original report only; no separate forcing theorem or precise skeleton definition is supplied.
\end{itemize}
\end{document}
